\section{Introduction}
Social conflicts increasingly generate parallel records in channels, comment threads, news reports, institutional statements, and short audiovisual fragments. These records support several different analytical tasks: documenting what a source said, tracing how a claim circulated, identifying a repeated text, organizing candidate events, and assessing whether available evidence supports an allegation. A system that conflates these tasks can turn a plausible narrative into an apparently verified account. A copied statement is not a new witness, a publication timestamp is not a recording timestamp, and a statistical departure from a timing model does not identify an organizer.

Crisis informatics has long emphasized the difficulty of extracting useful information from heterogeneous and noisy social-media streams \citep{imran2015crisis}. Computational social science adds concerns about platform access, sampling, measurement, and ethics \citep{Olteanu2019,VicentePaula2023}. Digital open-source investigation requires preservation and contextual verification of materials, including explicit treatment of uncertainty \citep{berkeleyProtocol2022}. These requirements become particularly important when software combines traditional analysis with generative AI. Fluent synthesis can hide uncertainty rather than resolve it; even citation-bearing generation requires separate assessment of whether citations support the text \citep{gao2023alce}.

Social addresses this integration problem as a research-software framework. Its practical objective is to help an analyst move from collected digital traces to inspectable intermediate records and bounded analytical outputs. The implementation contains platform-specific acquisition scripts, information-wave analysis, a linked evidence ledger, retrieval and verification rules, typed event processing, local media pipelines, and report validators. These components differ substantially in maturity, integration, and evaluation. Describing them as one uniformly autonomous investigator would obscure those differences.

This paper asks: \emph{How can a modular, evidence-grounded software system integrate heterogeneous digital traces to support automated analysis of social conflicts and information incidents while preserving provenance, uncertainty, and reproducibility?} We study an existing implementation rather than propose a new algorithm. The analysis freezes a Git revision, inspects execution paths, reconciles numerical artifacts, reruns bounded computations, and deliberately tests semantic failure cases. The Petropavlivska Borshchahivka application supplies realistic workloads, but is not an independent test of generalization. A separate project manuscript studies that incident itself \citep{pylypenko2026case}; the present paper examines the software and its evidentiary boundaries.

\subsection{Contributions of this work}
The contributions are fourfold. First, we provide a code-grounded integration map that separates reusable processors, platform collectors, experimental heuristics, curated case builders, and report-generation scripts. Second, we describe the implemented evidence model and the distinction between technical provenance, information lineage, source independence, and semantic validity. Third, we report reproducibility checks for similarity detection, conditional randomization, structural contracts, and a bounded offline workflow. Fourth, we expose concrete semantic counterexamples, measure a disclosed rule revision against them, and define the limits of the system's demonstrated autonomy. These contributions support more accountable use and evaluation of analytical software; they do not establish a new state of the art in fact verification, event extraction, or attribution.
